Critically, independent-sample $t$-tests demonstrated that the significant differences between diagnostic groups observed in real PET data were also observed in the synthetic images (Fig.~\ref{fig:4}C, D; figs.~\ref{fig:supp_disease_av45} and \ref{fig:supp_disease_fdg}). For instance, DIReCT$++$ recapitulated the statistically significant differences between AD and CN groups in the left precuneus for $^{18}$F-AV-45 (real: $q=7.73\times10^{-66}$; synthetic: $q=1.25\times10^{-63}$; Fig.~\ref{fig:4}C) and in the left hippocampus for $^{18}$F-FDG (real: $q=4.27\times10^{-15}$; synthetic: $q=1.27\times10^{-13}$; Fig.~\ref{fig:4}D). Similarly, regional differences that were non-significant in real PET data remained non-significant in the synthetic images after FDR correction (all $q\geq0.05$ for the regions shown in Fig.~\ref{fig:4}E, F). The AD--CN effect-size directions were consistent between real and synthetic PET for all ten regions shown in Fig.~\ref{fig:4}C, D, although the absolute Cohen's $d$ decreased in nine regions (maximum reduction, 0.0771). These findings support the preservation of group-level disease-related regional uptake patterns, while indicating variations in effect-size magnitude. Complete effect-size estimates and FDR-adjusted results for the displayed regions are provided in Supplementary Table~\ref{tab:regional_effect_sizes}. The AD--CN comparison illustrates representative differences between two diagnostic states rather than all diagnostic-group contrasts. MCI was evaluated separately through MCI-versus-CN and EMCI-versus-LMCI classification (Fig.~\ref{fig:5}B,C); the regional paired analysis also includes the original MCI categories.
